\textbf{Which method should one use?} The answer depends on what the decoder already holds. If nothing is shared in advance, the context-mixing compressors (zpaq, paq8px) give the smallest self-contained text files, at kilobytes per second; bsc and PPMd come within 10--15\% of them at megabytes per second; and xz, zstd and Brotli remain the practical choice when decoding speed matters. A pretrained language model is the smallest by a wide margin \emph{if} the decoder already has its weights and a GPU, which is plausible inside a system that runs the same model for other purposes, and implausible for archival or transmission between arbitrary machines. Our break-even sizes make this concrete: the knowledge in even the smallest useful model is repaid only after gigabytes of text.

\textbf{Multilingual fairness is mostly a tokeniser problem.} Haddadi and Teahan found a 1B-parameter LLaMA model far behind classical baselines on non-English text~\cite{haddadi2026comparing}, and Narashiman and Chandrachoodan saw GPT-2 expand Hindi~\cite{narashiman2024alphazip}. Our decomposition attributes the large majority of GPT-2's Hindi deficit, and all of its Chinese deficit, to the number of tokens its English-centred tokeniser produces, consistent with the tokenisation disparities measured by Petrov et al.~\cite{petrov2023tokenizers}. A multilingual tokeniser at a smaller model size (Qwen-0.5B) removes the deficit entirely, and the 1.5B model is then the best compressor we measured on Hindi. For compression, the choice of tokeniser matters more than a tripling of model size.

\textbf{Payload-only reporting misleads most for codebook methods.} EEHC's payload is small because each distinct block receives a short code, but the decoder needs every distinct block, so the codebook grows with the input's diversity. On natural-language text almost every 4-, 8- or 16-byte block is new, and the codebook makes the file larger than the input. The same bookkeeping applies, with different numbers, to DZip's stored bootstrap model and to language-model weights; reporting $r_P$, $r_F$ and $r_K$ side by side removes the ambiguity.

\textbf{Threats to validity.}\label{sec:threats}
(i) Slow methods ran on 1\,MB slices, which favours pretrained models and penalises online learners (TRACE, NNCP), which improve with more data; we state the slice size everywhere. (ii) NNCP ran in CPU mode with its default profile because its CUDA build crashed with our driver, and DZip and TRACE needed small compatibility fixes (listed in the repository); none of the fixes changes the coding. (iii) The largest language models used in prior work (7B and 70B parameters) do not fit our 6\,GB GPU; our models are 0.1--1.5B. (iv) Most runs shared the laptop with other jobs, so the speeds in \cref{fig:speed} are indicative; ratios are exact and every result was decoded and verified. (v) Our Python implementations of dynamic Huffman and rANS are not optimised; their speeds should not be compared with C implementations. (vi) The Leipzig sentence files are not parallel translations; cross-script per-character figures compare scripts and corpora, not languages.
